\documentclass[12pt,a4paper]{article}

% ── Packages ──────────────────────────────────────────────────────────────────
\usepackage[utf8]{inputenc}
\usepackage[T1]{fontenc}
\usepackage{lmodern}
\usepackage[british]{babel}
\usepackage[margin=2.5cm]{geometry}
\usepackage{setspace}
\usepackage{parskip}
\usepackage{enumitem}
\usepackage{booktabs}
\usepackage{longtable}
\usepackage{graphicx}
\usepackage{amsmath,amssymb}
\usepackage{xcolor}
\usepackage{csquotes}
\usepackage{hyperref}
\usepackage{array}
\usepackage{multirow}
\usepackage{tabularx}
\usepackage[style=apa,backend=biber,sortcites=true]{biblatex}
\addbibresource{bibliography.bib}

\hypersetup{
    colorlinks=true,
    linkcolor=blue!60!black,
    citecolor=blue!60!black,
    urlcolor=blue!60!black
}

\onehalfspacing

% ── Title ─────────────────────────────────────────────────────────────────────
\title{%
  \textbf{Data Collection Guide:\\
  Step-by-Step Instructions for GMT Revenue Simulation}%
}
\author{%
  Maurits de Haan\\
  \small MSc Economic Policy, Utrecht University\\
  \small Supervisor: dr.\ Peter Gerbrands%
}
\date{\today}

\begin{document}
\maketitle

\begin{abstract}
This appendix provides a step-by-step guide for collecting, downloading, and preparing the data necessary to conduct a simulation-based analysis of Global Minimum Tax (GMT) revenue effects.
Each data source is documented with its exact URL, download format, the variables it provides, the authors who have used it in the GMT literature, and the expected coverage in terms of jurisdictions, years, and MNE groups.
The guide is designed to ensure full replicability and to demonstrate the feasibility of the proposed research design.
\end{abstract}

\newpage
\tableofcontents
\newpage

% ══════════════════════════════════════════════════════════════════════════════
\section{Overview of the Data Architecture}\label{sec:overview}
% ══════════════════════════════════════════════════════════════════════════════

The proposed GMT revenue simulation requires data on five dimensions for each jurisdiction: (i)~pre-tax profits of multinational enterprises, (ii)~income taxes paid or accrued, (iii)~tangible assets, (iv)~number of employees and total payroll, and (v)~statutory and effective corporate tax rates.
Items~(i)--(iv) are needed to compute the GloBE effective tax rate and the substance-based income exclusion (SBIE); item~(v) is needed for calibration and cross-validation.
Additional auxiliary data---on Pillar~Two implementation status, country income classifications, and UPE locations---are required for the simulation of rule allocation and distributional analysis.

Table~\ref{tab:architecture} maps each data requirement to its primary source.

\begin{table}[!ht]
\centering
\caption{Data architecture: mapping research needs to sources}
\label{tab:architecture}
\small
\begin{tabular}{p{4cm}p{4.5cm}p{4cm}}
\toprule
\textbf{Data requirement} & \textbf{Primary source} & \textbf{Used by} \\
\midrule
MNE profits by jurisdiction & OECD CbCR (Table~I) & \textcite{barake2022}; \textcite{hugger2025}; \textcite{garcia2024} \\
\addlinespace
Income taxes paid/accrued & OECD CbCR (Table~I) & \textcite{barake2022}; \textcite{hugger2025} \\
\addlinespace
Tangible assets by jurisdiction & OECD CbCR (Table~I) & \textcite{devereux2023}; \textcite{barake2022} \\
\addlinespace
Employment \& payroll & OECD CbCR (Table~I) & \textcite{devereux2023}; \textcite{schjelderup2024} \\
\addlinespace
Statutory CIT rates & OECD Corporate Tax Statistics; Tax Foundation & \textcite{barake2022}; \textcite{torslov2022} \\
\addlinespace
Effective tax rates & OECD ETR database & \textcite{devereux2023}; \textcite{hugger2025} \\
\addlinespace
Pillar Two implementation status & OECD tracker; PwC tracker & \textcite{bilicka2023} \\
\addlinespace
UPE headquarters location & OECD CbCR (Table~III); Orbis & \textcite{barake2022} \\
\addlinespace
Country income classification & World Bank & \textcite{garcia2024} \\
\addlinespace
Profits \& wages of foreign firms & Eurostat FATS; BEA USDIA & \textcite{torslov2022} \\
\bottomrule
\end{tabular}
\end{table}


% ══════════════════════════════════════════════════════════════════════════════
\section{Step 1: OECD Country-by-Country Reporting Data}\label{sec:cbcr}
% ══════════════════════════════════════════════════════════════════════════════

The CbCR data are the \emph{central} dataset for any GMT simulation.
They were first released by the OECD in 2020 (covering fiscal year 2016) and are updated annually as part of the \emph{Corporate Tax Statistics} publication.
The data are anonymised and aggregated: individual MNE identities are suppressed, and values are summed across all MNE groups filing in a given ``reporting jurisdiction'' (i.e., the jurisdiction of the UPE or surrogate parent).

\subsection{What the data contain}

The CbCR dataset is organised into three tables:

\begin{description}[style=nextline,leftmargin=1.5cm]
\item[Table~I -- Aggregate totals by jurisdiction.]
For each \emph{reporting jurisdiction} (UPE location) and each \emph{partner jurisdiction} (where the MNE has entities), this table reports:
\begin{itemize}
    \item Number of MNE groups with entities in the partner jurisdiction
    \item Total number of constituent entities
    \item Total \textbf{unrelated party revenues}
    \item Total \textbf{related party revenues}
    \item Total \textbf{revenue} (= unrelated + related)
    \item \textbf{Profit (loss) before income tax}
    \item \textbf{Income tax paid (on cash basis)}
    \item \textbf{Income tax accrued -- current year}
    \item \textbf{Stated capital}
    \item \textbf{Accumulated earnings}
    \item \textbf{Number of employees}
    \item \textbf{Tangible assets other than cash and cash equivalents}
\end{itemize}

\item[Table~II -- Aggregate totals by sub-grouping.]
Similar to Table~I but disaggregated by the main business activity of the constituent entities (e.g., holding, manufacturing, services, licensing IP, etc.).

\item[Table~III -- Large MNE group composition (aggregated).]
Reports the number of MNE groups headquartered in each reporting jurisdiction that have entities in a given partner jurisdiction.
This table is essential for approximating UPE location for the IIR allocation.
\end{description}

\subsection{How to download}

\begin{enumerate}
\item Navigate to the OECD Corporate Tax Statistics database page:\\
\url{https://www.oecd.org/tax/beps/corporate-tax-statistics-database.htm}

\item Scroll to \textbf{``Anonymised \& aggregated Country-by-Country Reporting statistics''} under the ``Understanding the database'' section.

\item Click the direct link to the OECD Data Explorer:\\
\url{https://data-explorer.oecd.org/s/i6}

\item In the Data Explorer interface:
\begin{itemize}
    \item Select the desired \textbf{Table} (Table~I, Table~II, or Table~III) from the navigation panel.
    \item Select \textbf{Year}: all available years (currently 2016--2021, with 2022 expected by early 2026).
    \item Select \textbf{Reporting jurisdiction}: ``All'' to get the full cross-section.
    \item Select \textbf{Partner jurisdiction}: ``All''.
    \item Click \textbf{Download} $\rightarrow$ choose \textbf{CSV} format (recommended for large datasets) or \textbf{SDMX-ML} for programmatic access.
\end{itemize}

\item \textbf{Alternative: bulk download.} The full dataset is also available as a compressed CSV file via the OECD's SDMX REST API. The following URL downloads Table~I in TSV format:
\begin{quote}
\small\url{https://sdmx.oecd.org/public/rest/data/OECD.CTP.TPS,DSD_CBCR@DF_CBCR,1.0/all?format=csv}
\end{quote}
This can be imported directly into R or Python using \texttt{pandas.read\_csv()} or \texttt{readr::read\_csv()}.

\item \textbf{Companion documentation.} Download the explanatory annex and FAQ:
\begin{itemize}
    \item Explanatory Annex: \url{https://www.oecd.org/content/dam/oecd/en/topics/policy-sub-issues/corporate-taxation/explanatory-annex-corporate-tax-statistics.pdf}
    \item CbCR FAQ: \url{https://www.oecd.org/content/dam/oecd/en/topics/policy-sub-issues/corporate-taxation/corporate-tax-statistics-country-by-country-reporting-faqs.pdf}
\end{itemize}
These documents describe data limitations (double counting, aggregation thresholds, reporting gaps) and are essential reading before any analysis.
\end{enumerate}

\subsection{How authors have used this data}

\textcite{barake2022} use Table~I as the primary input for their GMT revenue simulation.
They compute the jurisdiction-level ETR for each reporting--partner jurisdiction pair as:
\[
\text{ETR}_{r,p} = \frac{\text{Income tax accrued (current year)}_{r,p}}{\text{Profit before income tax}_{r,p}}
\]
where $r$ denotes the reporting jurisdiction (UPE location) and $p$ denotes the partner jurisdiction where entities are located.
Pairs with negative profits are excluded.
The top-up tax for each pair below 15 per cent is then:
\[
\text{Top-up}_{r,p} = (\text{Profit}_{r,p} - \text{SBIE}_{r,p}) \times (0.15 - \text{ETR}_{r,p})
\]
where the SBIE is computed from the payroll and tangible asset data in the same table.

\textcite{hugger2025} access confidential \emph{micro-level} CbCR data at the OECD, which disaggregates values to the individual MNE--jurisdiction pair.
Their key finding---that 53.2 per cent of low-taxed profit is in jurisdictions with average ETRs above 15 per cent---cannot be directly replicated with public data but can be approximated using distributional assumptions (see Section~\ref{sec:coverage}).

\textcite{garcia2024} combine Table~I data with a logarithmic functional form for the profit-shifting relationship, using the CbCR profit and employee counts to identify ``excess'' profits in low-tax jurisdictions.

\subsection{Expected coverage}\label{sec:cbcr_coverage}

\begin{table}[!ht]
\centering
\caption{CbCR data coverage by fiscal year (as of November 2025)}
\label{tab:cbcr_coverage}
\small
\begin{tabular}{lccc}
\toprule
\textbf{Fiscal year} & \textbf{Reporting jurisdictions} & \textbf{Partner jurisdictions} & \textbf{MNE groups (approx.)} \\
\midrule
2016 & 26 & $\sim$160 & $\sim$4{,}000 \\
2017 & 38 & $\sim$175 & $\sim$6{,}000 \\
2018 & 47 & $\sim$190 & $\sim$7{,}000 \\
2019 & 49 & $\sim$195 & $\sim$7{,}500 \\
2020 & 50 & $\sim$200 & $\sim$7{,}500 \\
2021 & 52 & $\sim$200 & $\sim$8{,}000 \\
2022 & $\sim$55 (expected) & $\sim$200 & $\sim$8{,}500 \\
\bottomrule
\end{tabular}
\end{table}

\noindent\textbf{Key coverage notes:}
\begin{itemize}
    \item The \texteuro750~million revenue threshold means that only the largest MNEs are covered.
    These MNEs account for approximately 85--90 per cent of all profit shifting \parencite{oecd2025corptaxstats}.
    \item The \textbf{United States} is a reporting jurisdiction from 2016 onward, covering US-headquartered MNEs---the single most important group for profit-shifting analysis.
    \item \textbf{Tax havens} (Ireland, Luxembourg, Netherlands, Singapore, Switzerland) appear as partner jurisdictions in nearly all reporting jurisdictions' filings.
    Their coverage as reporting jurisdictions has improved over time (e.g., Ireland reports from 2017; Luxembourg from 2018).
    \item Partner jurisdictions are often grouped into residual categories (``Other Africa'', ``Other Asia'') when fewer than three MNE groups have entities there, to preserve anonymity.
\end{itemize}


% ══════════════════════════════════════════════════════════════════════════════
\section{Step 2: OECD Corporate Tax Statistics -- Tax Rates and ETRs}\label{sec:taxrates}
% ══════════════════════════════════════════════════════════════════════════════

\subsection{Statutory corporate income tax rates}

\begin{enumerate}
\item Navigate to the OECD Corporate Income Tax Rates Database:\\
\url{https://www.oecd.org/en/data/datasets/corporate-income-tax-rates-database.html}

\item Click \textbf{``Access the data''} or go directly to:\\
\url{https://data-explorer.oecd.org/s/hx}

\item Select:
\begin{itemize}
    \item \textbf{Rate type}: Combined (central + sub-central) corporate income tax rate.
    \item \textbf{Year}: All available years (1999--2025).
    \item \textbf{Country}: All Inclusive Framework jurisdictions.
\end{itemize}

\item Download as CSV.
The dataset covers approximately 145 Inclusive Framework jurisdictions, providing the statutory rate that determines the \emph{pre-GMT} tax environment.
\end{enumerate}

\subsection{Forward-looking effective tax rates}

The OECD also provides forward-looking ETRs computed using the Devereux--Griffith methodology:

\begin{enumerate}
\item Access via: \url{https://data-explorer.oecd.org/s/i0}
\item These ETRs are useful for cross-validation (comparing with backward-looking CbCR ETRs) and for understanding how the statutory system translates into effective burdens under different asset types and financing structures.
\end{enumerate}

\subsection{Tax Foundation data (supplementary)}

The Tax Foundation provides a comprehensive, annually updated dataset of statutory corporate tax rates for 225 jurisdictions:

\begin{enumerate}
\item Navigate to: \url{https://taxfoundation.org/data/all/global/corporate-tax-rates-by-country-2024/}
\item Download historical rates (1980--2024) as CSV:\\
\url{https://taxfoundation.org/wp-content/uploads/2025/01/rates_final.csv}
\item Download Pillar Two implementation data as Excel:\\
\url{https://taxfoundation.org/wp-content/uploads/2025/01/Statutory_Top_Corporate_Tax_Rates_and_Pillar_Two_Implementation.xlsx}
\item Download the underlying data from GitHub (version-controlled, machine-readable):\\
\url{https://github.com/TaxFoundation/worldwide-corporate-tax-rates}
\end{enumerate}

\textbf{How authors use this source.}
\textcite{barake2022} use the OECD statutory rates as a reference for identifying jurisdictions with rates below 15 per cent.
The Tax Foundation data complement this with broader country coverage and the critical information on which Pillar~Two rules each jurisdiction has adopted (IIR, QDMTT, UTPR).


% ══════════════════════════════════════════════════════════════════════════════
\section{Step 3: Pillar Two Implementation Tracker}\label{sec:tracker}
% ══════════════════════════════════════════════════════════════════════════════

The simulation of top-up tax allocation (QDMTT $\rightarrow$ IIR $\rightarrow$ UTPR) requires knowledge of which jurisdictions have implemented which rules and when.

\subsection{OECD Inclusive Framework tracker}

\begin{enumerate}
\item The OECD maintains an implementation tracker as part of the BEPS Inclusive Framework.
Access via the Corporate Tax Statistics page:\\
\url{https://www.oecd.org/tax/beps/corporate-tax-statistics-database.htm}
\item Scroll to \textbf{``BEPS Action 13 Implementation''} at: \url{https://data-explorer.oecd.org/s/i5}
\item The tracker records, for each Inclusive Framework jurisdiction, whether it has enacted legislation for the IIR, UTPR, and/or QDMTT, and the effective date.
\end{enumerate}

\subsection{PwC Pillar Two Country Tracker}

The most frequently updated and comprehensive tracker is maintained by PwC:
\begin{enumerate}
\item Navigate to:\\
\url{https://www.pwc.com/gx/en/services/tax/pillar-two-readiness/country-tracker.html}
\item The tracker provides, for each jurisdiction:
\begin{itemize}
    \item Whether the jurisdiction is an Inclusive Framework member
    \item Whether it has enacted IIR, UTPR, and/or QDMTT legislation
    \item The effective date of each rule
    \item Whether transitional safe harbours apply
    \item Links to enacted legislation
\end{itemize}
\item The data can be extracted manually into a spreadsheet or scraped programmatically.
As of early 2026, over 40 jurisdictions have enacted at least one GloBE rule.
\end{enumerate}

\subsection{How authors use this data}

\textcite{barake2022} use the implementation tracker to construct scenarios of partial adoption.
For each scenario (full adoption, EU-only, G7-only), they reassign top-up tax collections based on which rules are operational in which jurisdictions.
The tracker is essential for the IIR allocation: if the UPE jurisdiction has not enacted the IIR, the top-up tax remains uncollected (absent a UTPR backstop) or is allocated to a UTPR-implementing jurisdiction.


% ══════════════════════════════════════════════════════════════════════════════
\section{Step 4: Foreign Affiliates Statistics (FATS)}\label{sec:fats}
% ══════════════════════════════════════════════════════════════════════════════

FATS are a complementary data source used by \textcite{torslov2022} as the backbone of their macroeconomic approach to quantifying profit shifting.
They record the economic activity of foreign-controlled enterprises operating in a given host country.

\subsection{Eurostat Inward FATS}

\begin{enumerate}
\item Navigate to the Eurostat Data Browser:\\
\url{https://ec.europa.eu/eurostat/databrowser/explore/all/trade}
\item In the navigation tree, expand: \textbf{Industry, trade and services} $\rightarrow$ \textbf{Globalisation in business statistics (gbs)} $\rightarrow$ \textbf{Foreign controlled EU enterprises -- inward FATS (fats)}.
\item The key dataset is \texttt{fats\_activ}: ``Foreign controlled EU enterprises by NACE Rev.~2 activity and country of control.''\\
Direct link: \url{https://ec.europa.eu/eurostat/databrowser/product/view/fats_activ}
\item Variables include:
\begin{itemize}
    \item Number of enterprises
    \item Turnover
    \item \textbf{Value added at factor cost}
    \item \textbf{Personnel costs} (wages and salaries)
    \item \textbf{Gross operating surplus} (proxy for pre-tax profit)
    \item Number of persons employed
    \item Gross investment in tangible goods
\end{itemize}
\item Download in CSV or TSV format.
\item \textbf{Coverage:} All 27 EU member states plus Norway, Switzerland, and the UK (for years before 2021).
Data are available from 2008 (for selected countries) to 2022.
Breakdown by \textbf{country of control} allows identification of affiliates controlled by entities in specific jurisdictions (e.g., US-controlled, Irish-controlled).
\end{enumerate}

\subsection{BEA Activities of U.S.\ Multinational Enterprises}

For US-headquartered MNEs specifically, the Bureau of Economic Analysis (BEA) provides detailed data:

\begin{enumerate}
\item Navigate to: \url{https://www.bea.gov/international/di1usdop}
\item Under \textbf{``Frequently Requested Tables -- Majority-Owned Foreign Affiliates (MOFAs)''}:
\begin{itemize}
    \item \textbf{Selected Data by Country} (2009--2023):\\
    \url{https://www.bea.gov/sites/default/files/2025-08/selected-mofas-cntry.xls}
    \item \textbf{Employment by Country and Industry} (2009--2023):\\
    \url{https://www.bea.gov/sites/default/files/2025-08/mofas-employment.xls}
    \item \textbf{Assets by Country and Industry} (2009--2023):\\
    \url{https://www.bea.gov/sites/default/files/2025-08/mofas-assets.xls}
    \item \textbf{Value Added by Country and Industry} (2009--2023):\\
    \url{https://www.bea.gov/sites/default/files/2025-08/mofas-value-added.xls}
\end{itemize}
\item For more detailed (but older) data, use the latest Benchmark Survey (2019):\\
\url{https://www.bea.gov/worldwide-activities-us-multinational-enterprises-revised-2019-statistics}
\item BEA's interactive table system allows custom queries by country $\times$ industry:\\
\url{https://apps.bea.gov/iTable/iTableHtml.cfm?reqid=2&step=1&isuri=1}
\end{enumerate}

\subsection{How authors use FATS data}

\textcite{torslov2022} construct a profitability benchmark by comparing the \emph{profit-to-wage ratio} of foreign-controlled firms with that of local firms in the same jurisdiction.
Using Eurostat inward FATS, they observe that foreign firms in Ireland report a profit-to-wage ratio of 800 per cent, compared to 34 per cent for local firms---a discrepancy they attribute to profit shifting.
The FATS data avoid the double-counting problem that afflicts BEA data, because ``gross operating surplus'' excludes equity income by construction.

For the GMT simulation, FATS serve a supplementary role: they provide an independent estimate of the profit-to-wage and profit-to-asset ratios that can be used to calibrate the SBIE and to validate the CbCR-based ETR computations.


% ══════════════════════════════════════════════════════════════════════════════
\section{Step 5: World Bank Country Classifications}\label{sec:worldbank}
% ══════════════════════════════════════════════════════════════════════════════

The distributional analysis (research question C1) requires a classification of countries by income level.

\begin{enumerate}
\item Navigate to the World Bank Country and Lending Groups page:\\
\url{https://datahelpdesk.worldbank.org/knowledgebase/articles/906519-world-bank-country-and-lending-groups}
\item Download the current fiscal year classification (updated annually each July).
Alternatively, download the historical classification spreadsheet:\\
\url{https://datacatalog.worldbank.org/search/dataset/0037712/World-Development-Indicators}
\item The classification assigns each country to one of four groups:
\begin{itemize}
    \item Low income (GNI per capita $\leq$ \$1{,}145 in 2024)
    \item Lower-middle income (\$1{,}146--\$4{,}515)
    \item Upper-middle income (\$4{,}516--\$14{,}005)
    \item High income ($>$ \$14{,}005)
\end{itemize}
\item Map these classifications to the jurisdiction codes used in the CbCR data.
Note that the CbCR data use ISO 3166-1 alpha-3 codes, which correspond directly to World Bank codes.
\end{enumerate}

\textcite{garcia2024} use World Bank income classifications to show that lower-income countries lose a proportionally larger share of tax revenue to profit shifting.
The same classification allows the GMT simulation to report results by country income group, quantifying the distributional consequences of different rule-adoption scenarios.


% ══════════════════════════════════════════════════════════════════════════════
\section{Step 6: Supplementary Data for SBIE Calibration}\label{sec:sbie}
% ══════════════════════════════════════════════════════════════════════════════

The SBIE computation requires two variables at the jurisdiction level: eligible \textbf{payroll costs} and the carrying value of \textbf{tangible assets}.
Both are available in CbCR Table~I (``Number of employees'' and ``Tangible assets other than cash and cash equivalents'').
However, CbCR does not report payroll directly for all jurisdictions.

\subsection{Imputing payroll from employee counts}

Where payroll is not directly reported, it can be imputed using average wage data:

\begin{enumerate}
\item Download OECD Average Annual Wages:\\
\url{https://data-explorer.oecd.org/s/4h}
\item For non-OECD jurisdictions, use ILO Global Wage Report data:\\
\url{https://ilostat.ilo.org/topics/wages/}
\item Multiply: $\text{Payroll}_{r,p} \approx \text{Employees}_{r,p} \times \text{Average wage}_p$
\end{enumerate}

\textcite{devereux2023} follow this approach, using the CbCR employee count multiplied by national average wages to approximate the SBIE payroll component.
They document that the SBIE covers 57 per cent of pre-tax profits on average---a finding that can be validated with the downloaded data.

\subsection{SBIE formula}

The SBIE for each MNE--jurisdiction pair is computed as:
\[
\text{SBIE}_{r,p} = \alpha \times \text{Payroll}_{r,p} + \beta \times \text{Tangible assets}_{r,p}
\]
where $\alpha$ and $\beta$ are the exclusion percentages.
During the transitional period (2024--2033), these are:
\begin{itemize}
    \item 2024: $\alpha = 10\%$, $\beta = 8\%$
    \item 2025: $\alpha = 9.8\%$, $\beta = 7.8\%$
    \item \ldots (declining by 0.2 percentage points per year for payroll, varying for assets)
    \item 2033 onward: $\alpha = 5\%$, $\beta = 5\%$
\end{itemize}

The OECD GloBE Model Rules document (available at \url{https://www.oecd.org/tax/beps/tax-challenges-arising-from-the-digitalisation-of-the-economy-globe-anti-base-erosion-model-rules-pillar-two.htm}) provides the complete transitional schedule in Article~5.3.


% ══════════════════════════════════════════════════════════════════════════════
\section{Step 7: Data Integration and Preparation}\label{sec:integration}
% ══════════════════════════════════════════════════════════════════════════════

\subsection{Merging the datasets}

The following steps describe the data preparation pipeline:

\begin{enumerate}
\item \textbf{Load CbCR Table~I} into a data frame with columns: \texttt{reporting\_jur}, \texttt{partner\_jur}, \texttt{year}, \texttt{profit}, \texttt{tax\_accrued}, \texttt{tax\_paid}, \texttt{employees}, \texttt{tangible\_assets}, \texttt{revenue}.

\item \textbf{Compute ETR:}
\[
\texttt{etr} = \frac{\texttt{tax\_accrued}}{\texttt{profit}} \quad \text{(excluding rows where profit} \leq 0\text{)}
\]

\item \textbf{Merge statutory CIT rates} by \texttt{partner\_jur} and \texttt{year} from the OECD/Tax Foundation dataset.

\item \textbf{Merge average wages} by \texttt{partner\_jur} from OECD/ILO data.
Compute \texttt{payroll} = \texttt{employees} $\times$ \texttt{avg\_wage}.

\item \textbf{Compute SBIE:}
\[
\texttt{sbie} = \alpha \times \texttt{payroll} + \beta \times \texttt{tangible\_assets}
\]

\item \textbf{Compute top-up tax:}
\[
\texttt{topup} = \max\bigl(0,\; (\texttt{profit} - \texttt{sbie}) \times (0.15 - \texttt{etr})\bigr) \quad \text{if } \texttt{etr} < 0.15
\]

\item \textbf{Merge Pillar Two implementation tracker} by \texttt{partner\_jur} (for QDMTT status) and \texttt{reporting\_jur} (for IIR status).

\item \textbf{Allocate top-up tax} according to the rule order:
\begin{itemize}
    \item If the partner jurisdiction has enacted a QDMTT: allocate to the partner jurisdiction.
    \item Else, if the reporting jurisdiction (UPE location) has enacted the IIR: allocate to the reporting jurisdiction.
    \item Else, if any jurisdiction where the MNE operates has enacted the UTPR: allocate proportionally by tangible assets and employees in UTPR jurisdictions.
    \item Else: the top-up tax goes uncollected.
\end{itemize}

\item \textbf{Merge World Bank income classifications} by jurisdiction for distributional grouping.

\item \textbf{Aggregate results} by partner jurisdiction, reporting jurisdiction, income group, and scenario.
\end{enumerate}

\subsection{Software requirements}

The entire pipeline can be executed in \textbf{Python} (using \texttt{pandas}, \texttt{numpy}, and \texttt{matplotlib}) or \textbf{R} (using \texttt{tidyverse}, \texttt{readxl}, and \texttt{ggplot2}).
No proprietary software or restricted-access data are required.


% ══════════════════════════════════════════════════════════════════════════════
\section{Expected Data Coverage and Limitations}\label{sec:coverage}
% ══════════════════════════════════════════════════════════════════════════════

Table~\ref{tab:coverage} presents the expected coverage matrix showing which variables are available, for how many jurisdictions, and for which years.

\begin{table}[!ht]
\centering
\caption{Expected data coverage matrix}
\label{tab:coverage}
\small
\begin{tabular}{p{3.5cm}ccp{5cm}}
\toprule
\textbf{Variable} & \textbf{Jurisdictions} & \textbf{Years} & \textbf{Coverage notes} \\
\midrule
CbCR profits & $\sim$200 (as partner) & 2016--2021+ & Aggregated across MNE groups within each reporting jurisdiction; some partner jurisdictions grouped \\
\addlinespace
CbCR taxes paid & $\sim$200 & 2016--2021+ & Cash basis; may differ from accrued taxes \\
\addlinespace
CbCR taxes accrued & $\sim$200 & 2016--2021+ & ``Current year'' accruals; closer to GloBE covered taxes \\
\addlinespace
CbCR tangible assets & $\sim$200 & 2016--2021+ & ``Other than cash \& cash equivalents''; carrying value \\
\addlinespace
CbCR employees & $\sim$200 & 2016--2021+ & Full-time equivalent count \\
\addlinespace
Statutory CIT rates & $\sim$145 & 1999--2025 & Inclusive Framework jurisdictions; combined central + sub-central \\
\addlinespace
Forward-looking ETRs & $\sim$80 & 2017--2025 & OECD-calculated; based on Devereux--Griffith model \\
\addlinespace
Eurostat FATS & 30 (EU + EFTA) & 2008--2022 & Inward FATS: foreign-controlled firms by country of origin \\
\addlinespace
BEA MOFA data & $\sim$60 (host) & 2009--2023 & US-parented MNEs only; by country $\times$ industry \\
\addlinespace
Pillar Two tracker & $\sim$145 & 2024--2026 & Binary indicators for IIR/UTPR/QDMTT enactment dates \\
\addlinespace
World Bank income & 218 & Annual & Low / Lower-middle / Upper-middle / High \\
\bottomrule
\end{tabular}
\end{table}

\subsection{Visualising the coverage}

Figure~\ref{fig:coverage} illustrates the expected jurisdictional coverage of the CbCR data across the income distribution, showing that high-income countries are well covered as both reporting and partner jurisdictions, while low-income countries appear primarily as partner jurisdictions (and often in grouped residual categories).

\begin{figure}[!ht]
\centering
\fbox{%
\begin{minipage}{0.9\textwidth}
\centering
\vspace{0.3cm}
\textbf{Expected CbCR Coverage by Country Income Group}
\vspace{0.3cm}

\begin{tabular}{lcccc}
\toprule
& \textbf{High} & \textbf{Upper-mid} & \textbf{Lower-mid} & \textbf{Low} \\
\midrule
As reporting jurisdiction & 40+ & 8--10 & 2--3 & 0--1 \\
As partner jurisdiction (named) & 50+ & 30--40 & 20--30 & 5--15 \\
As partner (in residual group) & 0--2 & 10--15 & 15--25 & 20--30 \\
\addlinespace
\textbf{Share of global MNE profit covered} & $\sim$95\% & $\sim$3\% & $\sim$1.5\% & $<$0.5\% \\
\bottomrule
\end{tabular}
\vspace{0.3cm}
\end{minipage}
}
\caption{Expected CbCR coverage by World Bank income group.
High-income countries dominate both as UPE locations (reporting jurisdictions) and as identifiable partner jurisdictions.
Low-income countries are predominantly captured in residual categories, limiting jurisdiction-specific analysis.}
\label{fig:coverage}
\end{figure}


\subsection{Key limitations and mitigation strategies}

\begin{longtable}{p{4cm}p{4.5cm}p{4.5cm}}
\toprule
\textbf{Limitation} & \textbf{Impact on analysis} & \textbf{Mitigation} \\
\midrule
\endhead
\textbf{Aggregation across MNEs} -- CbCR data are summed across all MNE groups within each reporting--partner pair &
Cannot compute MNE-level ETRs; jurisdictional averaging conceals within-jurisdiction dispersion \parencite{hugger2025} &
Apply distributional assumptions (e.g., log-normal ETR distribution) calibrated to the Hugger et al.\ findings; conduct sensitivity analysis across distributional parameters \\
\addlinespace

\textbf{Double counting risk} -- Intercompany dividends and equity income may appear in multiple jurisdictions' profit figures &
Overstates total profit and understates ETRs &
Following \textcite{torslov2022}: use ``profit before tax'' which partially nets out equity income; apply proportional correction factors from the OECD FAQ document \\
\addlinespace

\textbf{CbCR vs.\ GloBE definitions} -- CbCR income and taxes follow national tax law / financial accounting, not GloBE-specific definitions &
ETR computed from CbCR may differ from the GloBE ETR &
This is the same limitation faced by the transitional CbCR safe harbour; document the discrepancy and quantify its likely magnitude using the safe harbour guidance \\
\addlinespace

\textbf{Residual jurisdiction groups} -- Small jurisdictions are merged into continental residual categories &
Some low-income countries cannot be individually identified &
Use FATS/BEA data and national sources to supplement where available; report results both with and without residual categories \\
\addlinespace

\textbf{Missing payroll data} -- CbCR reports employee count but not payroll directly &
SBIE payroll component must be imputed &
Impute using OECD/ILO average wage data, as done by \textcite{devereux2023}; conduct sensitivity analysis using wage ranges \\
\addlinespace

\textbf{Time lag} -- Most recent public CbCR data are for FY2021; FY2022 expected by early 2026 &
Analysis uses pre-implementation data only &
Clearly state that the analysis provides a pre-implementation baseline; compare with FY2022/2023 data if released in time \\
\bottomrule
\end{longtable}


% ══════════════════════════════════════════════════════════════════════════════
\section{Summary: Complete Download Checklist}\label{sec:checklist}
% ══════════════════════════════════════════════════════════════════════════════

Table~\ref{tab:checklist} provides a consolidated checklist of all data files to download, with exact URLs, formats, and estimated file sizes.

\begin{longtable}{p{0.4cm}p{3.5cm}p{5.5cm}p{1.2cm}p{1.2cm}}
\toprule
\textbf{\#} & \textbf{Dataset} & \textbf{URL} & \textbf{Format} & \textbf{Size} \\
\midrule
\endhead
1 & CbCR Table~I & \url{https://data-explorer.oecd.org/s/i6} & CSV & $\sim$50MB \\
\addlinespace
2 & CbCR Table~II & Same portal, select Table~II & CSV & $\sim$80MB \\
\addlinespace
3 & CbCR Table~III & Same portal, select Table~III & CSV & $\sim$15MB \\
\addlinespace
4 & Statutory CIT rates & \url{https://data-explorer.oecd.org/s/hx} & CSV & $\sim$1MB \\
\addlinespace
5 & Forward-looking ETRs & \url{https://data-explorer.oecd.org/s/i0} & CSV & $\sim$2MB \\
\addlinespace
6 & Tax Foundation rates & \url{https://github.com/TaxFoundation/worldwide-corporate-tax-rates} & CSV & $\sim$0.5MB \\
\addlinespace
7 & Tax Foundation Pillar~Two & \scriptsize\url{https://taxfoundation.org/wp-content/uploads/2025/01/Statutory_Top_Corporate_Tax_Rates_and_Pillar_Two_Implementation.xlsx} & XLSX & $\sim$0.3MB \\
\addlinespace
8 & PwC Pillar~Two tracker & \url{https://pwc.com/gx/en/services/tax/pillar-two-readiness/country-tracker.html} & Manual & --- \\
\addlinespace
9 & OECD Average Wages & \url{https://data-explorer.oecd.org/s/4h} & CSV & $\sim$0.5MB \\
\addlinespace
10 & ILO Wage data & \url{https://ilostat.ilo.org/topics/wages/} & CSV & $\sim$2MB \\
\addlinespace
11 & Eurostat inward FATS & \url{https://ec.europa.eu/eurostat/databrowser/product/view/fats_activ} & TSV & $\sim$30MB \\
\addlinespace
12 & BEA MOFA data & \url{https://www.bea.gov/international/di1usdop} & XLS & $\sim$5MB \\
\addlinespace
13 & World Bank income groups & \url{https://datahelpdesk.worldbank.org/knowledgebase/articles/906519} & XLS & $\sim$0.2MB \\
\addlinespace
14 & GloBE Model Rules & \scriptsize\url{https://www.oecd.org/tax/beps/tax-challenges-arising-from-the-digitalisation-of-the-economy-globe-anti-base-erosion-model-rules-pillar-two.htm} & PDF & --- \\
\addlinespace
15 & CbCR Explanatory Annex & \scriptsize\url{https://www.oecd.org/content/dam/oecd/en/topics/policy-sub-issues/corporate-taxation/explanatory-annex-corporate-tax-statistics.pdf} & PDF & --- \\
\bottomrule
\end{longtable}


% ══════════════════════════════════════════════════════════════════════════════
\section{Feasibility Assessment}\label{sec:feasibility}
% ══════════════════════════════════════════════════════════════════════════════

The data collection plan outlined above demonstrates that the proposed GMT revenue simulation is fully feasible using \textbf{publicly available, freely accessible} data.
The critical assessment is as follows:

\begin{description}[style=nextline,leftmargin=1.5cm]

\item[Core data availability: \textcolor{green!60!black}{CONFIRMED}.]
CbCR Table~I provides all five core variables (profits, taxes, tangible assets, employees, revenue) needed for the GloBE ETR computation and SBIE calculation.
These data are freely downloadable from the OECD Data Explorer.

\item[Jurisdiction coverage: \textcolor{green!60!black}{SUFFICIENT}.]
Approximately 200 jurisdictions appear as partner jurisdictions in Table~I, covering an estimated 85--90 per cent of global MNE profits.
The 50+ reporting jurisdictions include all major UPE locations (US, Japan, Germany, UK, France, Netherlands, China).

\item[Temporal coverage: \textcolor{orange!80!black}{ADEQUATE FOR BASELINE}.]
Data through FY2021 are currently available; FY2022 is expected by early 2026.
This provides 5--6 years of pre-implementation data for baseline construction.
Post-implementation data (FY2024+) will not be available within the thesis timeline.

\item[SBIE variables: \textcolor{green!60!black}{AVAILABLE WITH IMPUTATION}.]
Tangible assets are directly reported in CbCR.
Payroll must be imputed from employee counts and average wage data---a well-established approach in the literature \parencite{devereux2023,barake2022}.

\item[Rule allocation data: \textcolor{green!60!black}{AVAILABLE}.]
Pillar Two implementation status is tracked by both the OECD and PwC.
UPE location is identifiable from CbCR reporting jurisdictions and Table~III.

\item[Distributional data: \textcolor{green!60!black}{AVAILABLE}.]
World Bank income classifications are freely available and map directly to CbCR jurisdiction codes.

\item[Computational requirements: \textcolor{green!60!black}{MINIMAL}.]
All computations can be performed on a standard laptop using Python or R.
No proprietary software, high-performance computing, or restricted data access is required.

\end{description}

\vspace{0.5cm}
\noindent\fbox{%
\begin{minipage}{0.95\textwidth}
\textbf{Bottom line:} The proposed simulation-based analysis of GMT revenue effects---incorporating within-jurisdiction ETR dispersion approximation, SBIE sensitivity analysis, and distributional assessment across country income groups---can be executed entirely with freely available public data, standard software tools, and within the April--June 2026 thesis timeline.
The methodology replicates and extends the approaches used by \textcite{barake2022}, \textcite{devereux2023}, and \textcite{hugger2025}, while requiring no confidential data access.
\end{minipage}
}


% ══════════════════════════════════════════════════════════════════════════════
\printbibliography
% ══════════════════════════════════════════════════════════════════════════════

\end{document}
